\chapter{Conclusions}
We found that the organisation of the work matters at least as much as the number of threads. Task 1a gains substantially from partitioning a single file across workers: the primary time decreases from 3876 to 1277 ms, but returns beyond 8--16 threads are small and unstable between sessions. We interpret this as CPU-sensitive parsing combined with shared memory, I/O, scheduling, and runtime costs; without matched cold-cache tests or profiling, we do not describe the internal-disk task as purely CPU-bound.

We classify Task 1b as an in-memory CPU/memory-bound counting workload: input loading is outside the timed region, and the counting loops perform no file I/O. It is our clearest CPU-side comparison, although memory bandwidth and shared-cache traffic can still limit scaling. Per-thread counters are both correct and fastest in our eight-thread measurements. The locked version is correct but much slower because it serialises millions of small updates. The unsynchronised version returns different incorrect totals on the same input. Its elapsed time therefore cannot be used as evidence of a successful correct optimisation.

Task 1c demonstrates a different scaling limit. Its bounded queue prevents unlimited producer lead, but the single reader, string allocation, and queue transfers restrict throughput. The primary optimum occurs at two consumers, and 64 consumers are slower than one. Memory depends on queue occupancy and active batches as well as thread count, which explains why its peak RSS is not monotonic.

The cold USB experiment introduces a substantial storage-wait component: one thread needs \USBOneMs{} ms for the same dataset, compared with 3876 ms in the primary one-thread Task 1a session. The USB run uses only about 2.07 s of CPU time during 4.67 s of real time. Additional workers achieve the measured speed-ups of \USBSpeedTwo{} and \USBSpeedFour{}, while continuing to share one flash device. We therefore distinguish latency overlap from multiplying storage bandwidth, and we keep the USB conclusions separate from earlier measurements with unverified cache provenance.

Our practical recommendation for this workload is to use independent counters and a small, measured degree of concurrency, to verify all counters before interpreting timings, and to keep cache conditions explicit. The best worker count is an experimental result, not a value that can be inferred from core count alone.
